% ch2_a 流形 (书页 48-60 / p063-p075)
% 第2章 Manifolds
\chapter{流形}

% 2.1 Gravity as Geometry
\section{引力即几何}

引力是特殊的。在广义相对论的语境下，我们把这种特殊性归因于这样一个事实：产生引力的动力学场正是描写时空自身曲率的度规张量，而不是某种在时空中传播的附加场；这就是 Einstein 深刻的洞见。把他引向这一想法的物理原理，是引力相互作用的\emph{普适性}（universality），这一原理由\textbf{等效原理}（Principle of Equivalence）形式化。让我们看看这个物理原理如何把我们引向“将引力描述为弯曲流形（manifold）的几何”这一数学策略。

等效原理有各种不同的形式，其中第一种是\textbf{弱等效原理}（Weak Equivalence Principle），简称WEP。WEP断言：任何物体的惯性质量与引力质量相等。为了弄清这意味着什么，想想Newton力学。第二定律把作用在物体上的力与物体经历的加速度联系起来，令二者互成比例，比例常数就是惯性质量$m_i$：
\begin{equation*}
\mathbf{F} = m_i\mathbf{a}.
\tag{2.1}
\end{equation*}
惯性质量显然具有普适的特征：它与你试图推动物体时感受到的阻力有关，而且无论施加的是哪种力，它都取相同的值。我们还有Newton引力定律，它可以被理解为：作用在物体上的引力正比于某个标量场$\Phi$的梯度，这个$\Phi$称为引力势（gravitational potential）。此时比例常数称为引力质量$m_g$：
\begin{equation*}
\mathbf{F}_g = -m_g\nabla\Phi.
\tag{2.2}
\end{equation*}
表面看来，$m_g$与$m_i$的性质非常不同；它是引力所特有的一种量。如果你愿意，可以把$m_g/m_i$看成该物体的“引力荷”（gravitational charge）。尽管如此，Galileo早就证明（流传的说法是他从比萨斜塔上扔下重物，实际上他是让小球沿斜面滚下）：物质对引力的响应是普适的——在引力场中每个物体都以相同的速率下落，与物体的组成无关。在Newton力学中，这就翻译成WEP，它简单地就是
\begin{equation*}
m_i = m_g
\tag{2.3}
\end{equation*}
对任何物体都成立。一个直接的推论是，自由下落的检验粒子（test particle）的行为是普适的，与它们的质量（或它们可能具有的任何其他性质）无关；实际上，我们有
\begin{equation*}
\mathbf{a} = -\nabla\Phi.
\tag{2.4}
\end{equation*}
在实验上，下落物体的重力加速度与其组成无关这一点，已由Eötvös实验及其现代后继者验证到了极高的精度。

这提示了WEP的一个等价表述：存在一类穿过时空的特殊轨迹，称为\emph{惯性}（inertial）轨迹（或“自由下落”轨迹），不加速的粒子沿着它们运动——这里“不加速”的意思是“只受到引力作用”。显然，这对其他的力并不成立，比如电磁力。在电场中，带相反电荷的粒子将沿颇不相同的轨迹运动。与此相反，每个粒子都具有完全相同的引力荷。

WEP所蕴含的引力的普适性，还可以用另一种更流行的形式来表述。设想我们考虑一位待在密封得严严实实的箱子里的物理学家，她无法观察外面的世界，正在做涉及检验粒子运动的实验，比如测量局部引力场。当然，如果箱子是放在月球或木星上，她得到的答案会与放在地球上时不同。但如果箱子正以恒定的速率加速，答案也会不同；因为这会改变自由下落粒子相对于箱子的加速度。WEP意味着：仅仅通过观察自由下落粒子的行为，无法把引力场的效应与处于匀加速参考系中的效应区分开来。这是引力普适性的推论；相比之下，在电动力学中，通过观察带不同电荷的粒子的行为，是可以把匀加速与电磁场区分开来的。但对引力来说这不可能，因为“荷”必然正比于（惯性）质量。

谨慎起见，我们应该把“无法区分引力与匀加速”这一论断加以限制，只考虑“时空的足够小的区域”。如果密封箱足够大，引力场就会以可观察的方式随地点变化，而加速的效应却总是沿同一个方向。在火箭飞船或电梯里，粒子总是竖直下落。然而，在一个处于引力场中的非常大的箱子里，粒子会朝（比如说）地心运动，对于相距很远的实验来说，这就是不同的方向了。因此WEP可以表述如下：\emph{在时空的足够小的区域中，自由下落粒子在引力场中的运动与在匀加速参考系中的运动相同。}在更大的时空区域中，引力场会存在不均匀性，这会导致潮汐力（tidal force），而潮汐力是可以探测到的。

狭义相对论出现之后，质量这一概念失去了它的某些独特性，因为人们逐渐清楚，质量只不过是能量和动量（正如我们在第1章中看到的）的一种表现。因此，Einstein 很自然地想到把WEP推广成某种更具包容性的东西。他的想法很简单：无论箱中的物理学家做什么实验（不只是让检验粒子下落），她都绝对没有任何办法区分匀加速与外部引力场。这一合理的外推就成了如今所说的\textbf{Einstein 等效原理}（Einstein Equivalence Principle），简称EEP：\emph{在时空的足够小的区域中，物理定律归结为狭义相对论的定律；不可能借助局部实验探测引力场的存在。}

实际上，很难想象有什么理论遵守WEP却违反EEP。考虑一个氢原子，它是质子和电子的束缚态。它的质量其实小于质子与电子各自单独考虑时的质量之和，因为存在负的结合能——你必须把能量注入原子，才能把质子和电子分开。按照WEP，氢原子的引力质量因此小于其各组分质量之和；而引力场与电磁相互作用（正是它把原子束缚在一起）以恰到好处的方式耦合，使得引力质量恰好得出正确的值。这意味着，引力不仅必须普适地与静止质量耦合，还必须与所有形式的能量和动量耦合——这几乎就是EEP的论断。不过，反例也是可以构造的；比如，我们可以想象一种引力理论，在其中自由下落的粒子穿过引力场时开始旋转。这样，它们仍可以沿着与在加速参考系中相同的路径下落（从而满足WEP），但你仍然能够探测到引力场的存在（这就违反了EEP）。这类理论显得很做作，但自然界并没有什么定律禁止它们。

有时人们会在“引力的物理定律”与“非引力的物理定律”之间做出区分，而EEP被定义为只适用于后者。于是\textbf{强等效原理}（Strong Equivalence Principle，SEP）被定义为涵盖所有物理定律，无论是否与引力有关。一个违反SEP但不违反EEP的理论，是这样一个理论：其中的\emph{引力}结合能对物体惯性质量和引力质量的贡献并不相等；这样，举例来说，具有显著自引力的检验粒子（在这样一种概念说得通的范围内）就可能沿与较轻粒子不同的轨迹下落。

正是EEP蕴含（或至少暗示）了一点：我们应当把引力的作用归因于时空的曲率。记住，在狭义相对论（SR）中，惯性系扮演着突出的角色——虽然我们无法挑出某个参考系说它唯一地“静止”，却可以挑出一族“不加速的”（惯性）参考系。带电粒子在电磁场中的加速度因而可以相对这些参考系唯一地定义。而EEP则意味着引力是无法逃避的——根本不存在“引力中性物体”这种东西，使我们能相对它测量重力加速度。由此可知，重力加速度并不是某种能够被可靠定义的东西，因此也就没有多大用处。

相反，更有意义的做法是去\emph{定义}“不加速”就是“自由下落”，而我们也将这样做。由此我们被引向这样一个观念：引力不是一种“力”——力是导致加速度的东西，而我们关于零加速度的定义是“在碰巧存在于周围的无论什么引力场中自由地运动”。

这个看似无害的步骤对时空的本性有着深远的影响。在SR中，我们有一套程序：从某一点出发，通过把刚性杆连接起来并在杆上安装时钟，构造一个延伸至整个时空的惯性系。但由于引力场的不均匀性（又是因为这个原因），这不再可能。如果我们从某个自由下落的状态出发，用刚性杆搭建一个巨大的结构，那么在一段距离之外，自由下落的物体看起来就会相对这个参考系在加速，如图 2.1 所示。解决办法是保留惯性系的概念，但放弃它们能够被唯一地延伸至整个空间和时间的希望。作为替代，我们可以定义\textbf{局部惯性系}（locally inertial frames），即在时空足够小的区域内追随单个自由下落粒子运动的那些参考系。（每次我们说“足够小的区域”时，纯粹主义者应当想象一个极限过程，在其中我们把相应的时空体积取为零。）这是我们所能做到的最好程度，但它迫使我们放弃很多东西。例如，我们再也不能有把握地谈论遥远物体之间的相对速度，因为适合那些物体的惯性参考系与适合我们的完全不同。

作为物理学家，我们的工作是构造世界的数学模型，然后用观测和实验来检验这些模型的预言。追寻引力普适性的种种推论，已经使我们放弃了把引力表达为一种在时空中传播的力的想法，甚至放弃了参考系可以延伸至整个时空的全球性观念。

%FIG{2.1}: {全球性参考系的失效。由于每个粒子都感受到引力的影响，我们把"不加速"定义为"自由下落"。其后果是，无法再按照第1章概述的程序定义全球性的惯性坐标系，因为初始静止的粒子将开始相对这样一个参考系运动。}

%FIG{2.2}: {由相距为$z$、各自感受加速度$a$的两个火箭测得的多普勒频移。}

因此我们需要引入一种数学框架，使物理理论能够与这些结论相容。解决办法将是设想时空具有弯曲的几何，而引力是这种曲率的一种表现。用来描述曲率的恰当数学结构就是\emph{微分流形}（differentiable manifold）：本质上，它是一类局部看起来像平直空间、但整体几何可能非常不同的集合。（记住，EEP可以表述为“在时空的小区域中物理定律归结为狭义相对论的定律”，这与“局部类似于平直空间的集合”这一数学观念配合得很好。）

我们无法证明引力应当被看成时空的曲率；我们能做的是提出这个想法，推导它的种种后果，然后看结果是否与我们关于世界的经验合理地相符。让我们着手来做这件事。

考虑EEP最著名的预言之一：引力红移（gravitational redshift）。想象两个火箭相距为$z$，各自以某个恒定的加速度$a$在远离任何引力场的区域中运动，如图 2.2 所示。在时刻$t_0$，后方火箭发出一个波长为$\lambda_0$的光子。两火箭之间的距离保持不变，所以在我们的背景参考系中，光子经过时间$\Delta t = z/c$后到达前方火箭。（我们假设$\Delta v/c$很小，所以只算到一阶。）在这段时间里，火箭将获得附加速度$\Delta v = a\Delta t = az/c$。因此，到达前方火箭的光子将被通常的多普勒效应红移，红移量为
\begin{equation*}
\frac{\Delta\lambda}{\lambda_0} = \frac{\Delta v}{c} = \frac{az}{c^2}.
\tag{2.5}
\end{equation*}
按照EEP，同样的事情也应该发生在一个均匀引力场中。于是我们想象一座高度为$z$的塔矗立在某颗行星的表面上，引力场强度为$a_g$（也就是Newton所说的“重力加速度”），如图 2.3 所示。我们设想塔顶火箭中的观者能够探测到从地面发出的光子，但除此之外他们无法看到外面，从而不知道自己正坐在一座塔顶。换句话说，他们没有办法把这种情形与加速火箭的情形区分开。因此，EEP让我们能够立即得出结论：从地面发出的波长为$\lambda_0$的光子将被红移，红移量为
\begin{equation*}
\frac{\Delta\lambda}{\lambda_0} = \frac{a_g z}{c^2}.
\tag{2.6}
\end{equation*}
这就是著名的引力红移。注意，它是EEP的直接推论；并不需要广义相对论的细节。

%FIG{2.3}: {地球表面上的引力红移，由不同高度的观者测量。}

红移公式更常用Newton势$\Phi$来表述，其中$\mathbf{a}_g = \nabla\Phi$。（符号与通常的约定相反，因为我们这里把$\mathbf{a}_g$看成参考系的加速度，而不是粒子相对该参考系的加速度。）$\Phi$的非常数梯度类似于一个随时间变化的加速度，等价的净速度由对光子从发射到吸收之间的时间积分给出。于是我们有
\begin{align*}
\frac{\Delta\lambda}{\lambda_0} &= \frac{1}{c}\int \nabla\Phi\, dt\\
&= \frac{1}{c^2}\int \partial_z\Phi\, dz\\
&= \Delta\Phi,
\tag{2.7}
\end{align*}
其中$\Delta\Phi$是引力势的总变化，而且我们再一次令$c=1$。这个简单的引力红移公式在更普遍的情形下依然成立。当然，只要用到了Newton势，我们就把适用范围限制在了弱引力场。

我们已经从EEP出发论证了引力红移的存在；现在可以用这个现象来进一步支持“应当把时空看成弯曲的”这一观念。考虑与之前相同的实验装置，只是现在画在图 2.4 的时空图中。地面上的物理学家从高度$z_0$处发出一束波长为$\lambda_0$的光，光传播到高度为$z_1$的塔顶。从光的任意一个波长的起点被发出，到这同一个波长的终点被发出，二者之间的时间间隔是$\Delta t_0 = \lambda_0/c$；而吸收端相应的时间间隔是$\Delta t_1 = \lambda_1/c$，时间分别由位于两个高度处的时钟测量。既然我们设想引力场是静态的，单个波的波前与波尾所经过的时空路径就必须严格全等。（图中用一般的弯曲路径来表示它们，因为我们并不假装知道这些路径究竟是什么样子。）简单的几何似乎意味着$\Delta t_0$与$\Delta t_1$必定相同。但它们当然不相同；引力红移意味着高处的实验者每秒观测到的波长数目更少，所以$\Delta t_1 > \Delta t_0$。我们可以粗略地把这解释为“塔上的钟似乎走得更快”。那么，哪里出了问题？出在简单的几何上——光子所穿过的时空是弯曲的。

%FIG{2.4}: {图 2.3 所绘引力红移实验的时空图。不同时刻出发的时空路径是全等的，但地面与塔上测得的时间间隔不同，这标志着欧几里得几何的失效。}

因此，我们希望把时空描述成这样一种数学结构：它在局部看起来像 Minkowski 空间，但在延展的区域上可能具有非平凡的曲率。涵盖这一观念的那类对象就是流形。在本章中，我们将只限于理解流形的概念以及可以在其上定义的各种结构，而把曲率的精确刻画留到下一章。

% 2.2 What Is a Manifold?
\section{什么是流形？}

流形（或微分流形）是数学和物理学中最基本的概念之一。我们都熟悉$n$维欧几里得空间$\mathbf{R}^n$的性质，它是$n$元组$(x^1,\ldots,x^n)$的集合，通常还配备一个平直的正定度规，其分量为$\delta_{ij}$。数学家们多年来致力于发展$\mathbf{R}^n$中的分析理论——微分、积分、函数的性质等等。但显然还存在其他一些空间（比如球面），我们直觉上认为它们是“弯曲的”，或者也许在拓扑上很复杂，我们也希望在它们上面进行类似的操作。

为了解决这个问题，我们发明了流形的概念，它对应于这样一种空间：可以是弯曲的，可以有复杂的拓扑，但在局部区域看起来就像$\mathbf{R}^n$。这里“看起来像”并不是说度规相同，而只是说一些更原初的概念——比如函数和坐标——以类似的方式运作。整个流形是通过把这些局部区域光滑地缝合在一起而构造出来的。一个关键的要点是，所使用的欧几里得空间的维数$n$在流形的每一片（patch）上都必须相同；这时我们就说这个流形的维数是$n$。用这种方法，我们可以通过把函数（局部地）转化为欧几里得空间中的函数，来分析这类空间上的函数。流形的例子包括：

\begin{itemize}
\item $\mathbf{R}^n$本身，包括直线（$\mathbf{R}$）、平面（$\mathbf{R}^2$）等等。这一点应当是显然的，因为$\mathbf{R}^n$不仅在局部而且在整体上都看起来像$\mathbf{R}^n$。
\item $n$维球面（$n$-sphere）$S^n$。它可以定义为$\mathbf{R}^{n+1}$中与原点相距某个固定距离的所有点的轨迹。圆周当然是$S^1$，而二维球面$S^2$是流形最有用的例子之一。（零维球面$S^0$，如果你仔细想想，由两个点组成。我们说$S^0$是一个不连通的零维流形。）值得强调的是，借助在$\mathbf{R}^{n+1}$中的嵌入来定义$S^n$只是一个方便的捷径；我们将讨论的所有流形都可以就其自身来定义，而无需求助于更高维的平直空间。
\item $n$维环面（$n$-torus）$T^n$，由取一个$n$维立方体并认同其相对的面而得到。二维环面$T^2$是一个认同了对边的正方形，如图 2.5 所示。甜甜圈的表面是一个大家熟悉的例子。
\item 亏格（genus）为$g$的黎曼面（Riemann surface），本质上是有$g$个洞而不止一个洞的二维环面，如图 2.6 所示。$S^2$可以看成亏格为零的黎曼面。用专业术语来说（这与我们眼下的讨论关系不大），每一个“紧致、可定向、无边界”的二维流形都是某个亏格的黎曼面。
\item 更抽象一些，由连续变换构成的集合，比如$\mathbf{R}^n$中的转动，构成一个流形。李群（Lie group）是同时还具有群结构的流形。例如SO(2)，即二维空间中转动的集合，就与$S^1$是同一个流形（尽管一般而言群流形会比球面更复杂）。
\item 两个流形的直积（direct product）是流形。也就是说，给定维数为$n$和$n'$的流形$M$和$M'$，我们可以构造一个维数为$n+n'$的流形$M\times M'$，它由有序对$(p,p')$组成，其中$p\in M$，$p'\in M'$。
\end{itemize}

%FIG{2.5}: {通过认同正方形的对边构造出的环面$T^2$。}

%FIG{2.6}: {不同亏格的黎曼面（genera是"genus"的复数）。}

看了这么多例子，流形的概念可能显得空洞无物：哪些东西\emph{不是}流形呢？很多东西都不是流形，因为它们在某些地方局部看起来不像$\mathbf{R}^n$。例子包括：一条撞进二维平面里的一维直线，以及在顶点处粘在一起的两个锥面，如图 2.7 所示。更微妙的例子见图 2.8。比如考虑一个单独的（二维）锥面。显然，在某种意义上锥面局部看起来像$\mathbf{R}^2$；与此同时，同样明显的是，锥面的顶点处有某种奇异的东西。这正是“微分流形”中“可微”一词开始发挥作用的地方；正如我们在发展形式定义时将会看到的，作为流形，锥面是完全光滑的，尽管其曲率在顶点处并不光滑。（其他类型的奇异性更加严重，会使我们无法把某些空间看成流形——无论光滑与否。）另一个例子是线段（包括端点）。由于端点的存在，它肯定不符合我们将要发展的流形定义。尽管如此，我们可以把定义加以推广，把“带边流形”（manifolds with boundary）包括进来，线段就是它们的范例。关于带边流形的简要讨论见附录D。

%FIG{2.7}: {不是流形的空间的例子：一条终止于平面上的直线，以及在顶点处相交的两个锥面。每种情形中都存在一个点，它在局部看起来不像固定维数的欧几里得空间。}

%FIG{2.8}: {微妙的例子。单个锥面是光滑流形，尽管其曲率在顶点处奇异。线段不是流形，但可以用更一般的"带边流形"概念来描述。}

这些微妙的情形应该会让你相信，一个严格的定义是必要的，我们现在就开始构造它；这里的讨论遵循Wald (1984)的路数。流形的非正式想法是：一个由一些局部看起来像$\mathbf{R}^n$的“片”（patch）光滑地缝合而成的空间。因此我们需要把“局部看起来像$\mathbf{R}^n$”和“光滑地缝合在一起”这两个观念形式化。我们需要若干预备定义，其中大多数都相当明白，但完整一点总是好的。最基本的概念是两个集合之间的\textbf{映射}（map）。（我们假定你知道集合是什么，或者你认为你知道；我们不需要太精确。）给定两个集合$M$和$N$，映射$\phi:M\to N$是这样一种关系：它把$M$的每个元素恰好指派为$N$的一个元素。因此，映射不过是函数的简单推广。给定两个映射$\phi:A\to B$和$\psi:B\to C$，我们通过运算$(\psi\circ\phi)(a)=\psi(\phi(a))$定义\textbf{复合}（composition）$\psi\circ\phi:A\to C$，如图 2.9 所示。于是$a\in A$，$\phi(a)\in B$，从而$(\psi\circ\phi)(a)\in C$。映射书写的次序是有意义的，因为右边的先作用。

如果$N$的每个元素至多有一个$M$的元素被映入其中，就称映射$\phi$是\textbf{一一的}（one-to-one，或单射/injective）；如果$N$的每个元素至少有一个$M$的元素被映入其中，就称它是\textbf{映上的}（onto，或满射/surjective）。（想一想就会发现，“one-to-one”更好的名字应该是“one-from-one”，或者干脆叫“two-to-two”。）考虑函数$\phi:\mathbf{R}\to\mathbf{R}$。那么$\phi(x)=e^x$是一一的但不是映上的；$\phi(x)=x^3-x$是映上的但不是一一的；$\phi(x)=x^3$两者都是；而$\phi(x)=x^2$两者都不是，如图 2.10 所示。

%FIG{2.9}: {映射$\psi\circ\phi:A\to C$由$\phi:A\to B$与$\psi:B\to C$复合而成。}

%FIG{2.10}: {几类映射。}

集合$M$称为映射$\phi$的\textbf{定义域}（domain），$M$被映入的$N$中那些点的集合称为$\phi$的\textbf{像}（image）。对任何子集$U\subset N$，$M$中被映到$U$的元素的集合称为$U$在$\phi$下的\textbf{原像}（preimage），记作$\phi^{-1}(U)$。既一一又映上的映射称为\textbf{可逆的}（invertible，或双射/bijective）。这时我们可以定义\textbf{逆映射}（inverse map）$\phi^{-1}:N\to M$，它满足$(\phi^{-1}\circ\phi)(a)=a$，如图 2.11 所示。注意，同一个符号$\phi^{-1}$既用于原像也用于逆映射，尽管前者总是有定义，而后者只在某些特殊情形下才有定义。

%FIG{2.11}: {一个映射及其逆映射。}

映射的\textbf{连续性}（continuity）概念实际上是一个非常微妙的概念，其精确表述我们并不需要。我们将假定你理解连续性和可微性这两个概念对于通常的函数（即映射$\phi:\mathbf{R}\to\mathbf{R}$）的含义。接下来，把这些概念推广到更一般的欧几里得空间之间的映射$\phi:\mathbf{R}^m\to\mathbf{R}^n$将会很有用。从$\mathbf{R}^m$到$\mathbf{R}^n$的映射把一个$m$元组$(x^1,x^2,\ldots,x^m)$变成一个$n$元组$(y^1,y^2,\ldots,y^n)$，因此可以看成$n$个$m$元函数$\phi^i$的汇集：
\begin{equation*}
\begin{gathered}
y^1 = \phi^1(x^1, x^2, \ldots, x^m)\\
y^2 = \phi^2(x^1, x^2, \ldots, x^m)\\
y^n = \phi^n(x^1, x^2, \ldots, x^m).
\end{gathered}
\tag{2.8}
\end{equation*}
如果这些函数中的任何一个的$p$阶导数存在且连续，我们就称它为$C^p$的；如果$\phi:\mathbf{R}^m\to\mathbf{R}^n$的每个分量函数都至少是$C^p$的，我们就称整个映射为$C^p$的。于是，$C^0$映射连续但不一定可微，而$C^\infty$映射连续并且可以想微分多少次就微分多少次。例如考虑一元函数$\phi(x)=|x^3|$。这个函数除$x=0$外处处无穷可微，而在$x=0$处可微两次但不能三次；因此我们说它是$C^2$的。$C^\infty$映射有时称为\textbf{光滑}（smooth）。

如果存在一个$C^\infty$映射$\phi:M\to N$，其逆$\phi^{-1}:N\to M$也是$C^\infty$的，我们就称两个集合$M$和$N$是\textbf{微分同胚的}（diffeomorphic）；这时映射$\phi$就称为一个\textbf{微分同胚}（diffeomorphism）。这是我们关于“两个空间作为流形是相同的”所拥有的最好概念。例如，当我们说SO(2)与$S^1$是同一个流形时，意思就是它们微分同胚。更多的讨论见附录B。

这些基本定义你可能早已熟悉，哪怕只是记得个大概。现在我们将把它们用于流形的严格定义。遗憾的是，要把这个相对直观的概念形式化，需要一套多少有些繁复的程序。我们必须首先定义开集的概念，在开集上我们可以放置坐标系，然后再以恰当的方式把开集缝合起来。

我们从\textbf{开球}（open ball）的概念开始，它是$\mathbf{R}^n$中所有满足$|x-y|<r$的点$x$的集合，其中$y\in\mathbf{R}^n$和$r\in\mathbf{R}$固定，而$|x-y|=\left[\sum_i(x^i-y^i)^2\right]^{1/2}$。注意这是一个严格不等式——开球就是以$y$为中心、半径为$r$的$n$维球面的内部，如图 2.12 所示。$\mathbf{R}^n$中的\textbf{开集}（open set）是由开球的任意（可能是无限的）并构造出来的集合。换句话说，$V\subset\mathbf{R}^n$是开的，如果对任何$y\in V$，都存在一个以$y$为中心、完全位于$V$内部的开球。粗略地说，开集就是某个$(n-1)$维闭曲面的内部（或者若干个这种内部的并）。

%FIG{2.12}: {定义在$\mathbf{R}^n$中的开球。}

\textbf{卡片}（chart）或\textbf{坐标系}（coordinate system）由一个集合$M$的子集$U$连同某个一一映射$\phi:U\to\mathbf{R}^n$组成，且像$\phi(U)$在$\mathbf{R}^n$中是开的，如图 2.13 所示。（任何映射都是映到其像上的，所以只要$\phi:U\to\phi(U)$是一一的，它就可逆。）这样我们就可以说$U$是$M$中的一个开集。一个$C^\infty$的\textbf{图册}（atlas）是卡片的带指标的汇集$\{(U_\alpha,\phi_\alpha)\}$，它满足两个条件：
\begin{enumerate}
\item 各$U_\alpha$的并等于$M$；也就是说，诸$U_\alpha$覆盖$M$。
\item 各卡片光滑地缝合在一起。更精确地说，如果两张卡片重叠，$U_\alpha\cap U_\beta\neq\varnothing$，那么映射$(\phi_\alpha\circ\phi_\beta^{-1})$把$\phi_\beta(U_\alpha\cap U_\beta)\subset\mathbf{R}^n$中的点\emph{映上}到开集$\phi_\alpha(U_\alpha\cap U_\beta)\subset\mathbf{R}^n$，并且所有这些映射在其有定义之处必须是$C^\infty$的。参照图 2.14（改编自Wald (1984)）应当会更清楚。
\end{enumerate}
所以，卡片就是我们通常所想的某个开集上的坐标系，而图册就是在重叠处彼此光滑相容的一个卡片系统。

%FIG{2.13}: {覆盖$M$的开子集$U$的一个坐标卡片。}

%FIG{2.14}: {相互重叠的坐标卡片。}

于是，经过漫长的铺垫：一个$C^\infty$的$n$维\textbf{流形}（manifold，或简称$n$维流形）不过是一个集合$M$连同它的一个极大图册，即包含每一个可能的相容卡片的图册。在上述所有定义中，我们也可以把$C^\infty$换成$C^p$。对我们的目的来说，流形的可微程度并不关键；我们总是假定任何流形都具有当前应用所必需的可微性。要求图册是极大的，是为了让配备了不同图册的两个等价空间不被算作不同的流形。这个定义用形式化的语言捕捉了我们关于“局部看起来像$\mathbf{R}^n$的集合”的观念。当然，我们很少会需要用上这个定义的全部力量，但精确本身就是一种回报。

我们的这个定义有一个好处：它不依赖于把流形嵌入某个更高维的欧几里得空间。实际上，任何$n$维流形都可以嵌入$\mathbf{R}^{2n}$中（Whitney嵌入定理），有时我们也会用到这个事实，比如上面我们对球面的定义。（克莱因瓶就是一个二维流形的例子，它不能嵌入$\mathbf{R}^3$，尽管可以嵌入$\mathbf{R}^4$。）但重要的是要认识到，流形具有独立于任何嵌入的单独存在性。例如，我们没有必要相信四维时空是被塞在某个更大的空间里的。不过话说回来，它也可能是；我们真的不知道。弦论的近期进展使人提出这样的猜想：我们可见的宇宙其实也许是更高维空间中的一张“膜”（brane，“membrane”的推广）。但就经典GR而言，四维的图景完全够用。

为什么有必要在卡片及其重叠的问题上如此讲究，而不是干脆用一张卡片覆盖每一个流形？因为大多数流形无法只用一张卡片覆盖。考虑最简单的例子$S^1$。有一个常规的坐标系$\theta:S^1\to\mathbf{R}$，其中$\theta=0$取在圆周的顶部，绕一圈直到$2\pi$。然而，在卡片的定义中我们已经要求像$\theta(S^1)$在$\mathbf{R}$中是开的。如果取$\theta=0$或$\theta=2\pi$中的任何一个，我们得到的都是一个闭区间而不是开区间；如果两个点都不取，我们
% 本文件至此为止：书页60末句止于"如果两个点都不取，我们"，正文延续到书页61，由下一块接续。
